% Definitions and timing. No payoff or optimization derivation.
\section{Primitives, Objects, and Timing}
\label{sec:p01-primitives-timing}

This section fixes the objects and information structure of the baseline before
any payoff is constructed.  The model is a partial equilibrium of
original-drug project advancement, commercialization organization, and
qualified manufacturing services.  The MAH reform changes only the
institutional feasibility or friction of retained entrusted manufacturing.

\subsection{Units and populations}
\label{subsec:p01-units-populations}

Let \(\mathsf C\) denote currency units, \(\mathsf X\) project-advancement
input units, \(\mathsf P\) the expected measure of viable planning-stage
projects in one decision cohort, \(\mathsf K\) manufacturing
capability/requirement units, and \(\mathsf B\) qualified manufacturing-service
capacity units.

There is a continuum \(\mathcal I\) of original-drug developers.  Developer
characteristics follow the exogenous joint distribution \(H(a,k)\), with
no independence restriction. Developer \(i\)
has project-advancement capability \(a_i>0\), measured in \(\mathsf P/\mathsf X\), and
internal manufacturing capability \(k_i>0\), measured in \(\mathsf K\). It
has the baseline characteristic vector
\begin{equation}
  \theta_i=(a_i,k_i).
  \label{eq:p01-developer-type}
\end{equation}
These are continuous characteristics, not permanent discrete firm types.

There is also a set \(\mathcal J\) of qualified manufacturing-service
suppliers.  Supplier efficiency \(z_j>0\) follows the exogenous distribution
\(H_C(z)\).  The supplier's capacity control and cost function are introduced
when the qualified manufacturing-service market is defined below.

\subsection{Project characteristics and empirical classification}
\label{subsec:p01-projects}

A viable project that reaches the route-planning stage draws commercial or
scientific value shifter \(q>0\) and manufacturing requirement \(m>0\).  The
latter is measured on the same \(\mathsf K\) scale as \(k_i\).  Define
\begin{equation}
  \omega=(q,m)\sim F(q,m),
  \label{eq:p01-project-draw}
\end{equation}
where \(F\) is an exogenous joint distribution.  Thus \(q\) governs the
project's value side, while the pair \((m,k_i)\) governs the manufacturing
mismatch.  The institutional regime does not directly alter \(q\), \(m\), or
\(F\).

For empirical decompositions only, let
\(g\in\{O,\mathrm{Inc}\}\) be an exogenous novelty classifier.  The label
\(\mathrm{Inc}\) cannot be confused with internal route \(I\).  Let
\(\rho_g\geq0\) be the class share and \(F_g(q,m)\) the conditional project
distribution.  They satisfy
\begin{equation}
  \sum_{g\in\{O,\mathrm{Inc}\}}\rho_g=1,
  \qquad
  F(q,m)=
  \sum_{g\in\{O,\mathrm{Inc}\}}\rho_gF_g(q,m).
  \label{eq:p01-class-mixture}
\end{equation}
The classifier is not a new baseline state, and it does not create a separate
control \(x_{ig}\).

\subsection{Institutional regime and route labels}
\label{subsec:p01-institution-routes}

Let the exogenous institutional regime be \(M\in\{0,1\}\).  The value \(M=0\)
denotes the pre-MAH regime.  The value \(M=1\) means that a developer may
legally retain authorization while a qualified external producer manufactures
the product.  Let \(\tau_E(M)\), measured in \(\mathsf C\) per project, denote
the institutional barrier attached only to this retained entrusted route.  The
baseline representation is
\begin{equation}
  \tau_E(0)=+\infty,
  \qquad
  \tau_E(1)=\bar\tau_E<+\infty,
  \qquad
  \bar\tau_E\geq0.
  \label{eq:p01-institutional-wedge}
\end{equation}
There is no continuous implementation parameter in the baseline.

The route labels are \(I\) for internal manufacturing, \(E\) for retained
entrusted manufacturing, \(T\) for transfer or another non-retained outside
option, and \(A\) for abandonment or indefinite delay.  If \(r_i\) denotes the
project's route control, its common label domain is
\begin{equation}
  r_i\in\{I,E,T,A\}.
  \label{eq:p01-route-domain}
\end{equation}
Effective pre-MAH unavailability of \(E\) is encoded solely by
\(\tau_E(0)=+\infty\), which gives \(W_i^E=-\infty\) once route values are
assembled.  Thus the common label domain is not a second policy channel.
Under \(E\), the developer remains the holder; entrusted production is not a
transfer of ownership or authorization.  The eventual optimizer is denoted by
\(r_i^*(q,m;M,p_m)\), but route values and the argmax are defined in the
organizational-sorting section below.

\subsection{Project-advancement and realization objects}
\label{subsec:p01-advancement-realization}

Developer \(i\) chooses one common control \(x_i\geq0\), measured in
\(\mathsf X\).  Canonically, \(x_i\) is
\emph{original-drug innovation investment / project-advancement intensity}.
It includes investment that advances viable original-drug projects toward a
commercialization-relevant route-planning stage.  It is broader than pure
clinical-development effort, but it is not basic research, basic-compound
discovery, patent-generating effort, or patent applications.

Project-advancement capability converts the control into the expected measure of viable
planning-stage projects:
\begin{equation}
  \lambda_i^{\mathrm{plan}}=a_ix_i.
  \label{eq:p01-planned-intensity}
\end{equation}
The control \(x_i\) has units \(\mathsf X\), so
\(\lambda_i^{\mathrm{plan}}\) has units \(\mathsf P\).

Let \(p_m\), measured in \(\mathsf C/\mathsf B\), denote a conjectured price
of qualified manufacturing service and let \(p_m^*\) denote its future
market-consistent value.  Let \(\Omega_i(M,p_m)\), measured in \(\mathsf C\)
per planning-stage project, denote developer \(i\)'s expected optimized project
value.  These objects are semantically reserved here: the sections below first
build route values, then solve the advancement problem, and finally close the
CMO fixed point.

Finally, let \(s(q)\in[0,1]\) be the exogenous, route-independent probability
of downstream clinical or regulatory realization:
\begin{equation}
  \Pr\{\text{downstream realization}\mid q\}=s(q).
  \label{eq:p01-downstream-probability}
\end{equation}
The reform does not shift \(s(q)\).  If typed observed outcomes later require a
class-specific probability \(s_g(q)\), it is likewise exogenous and
MAH-invariant.

\subsection{Baseline timing}
\label{subsec:p01-timing}

\begin{enumerate}
  \item[Stage 0.] \textbf{Institutional regime.}  The value of \(M\) is
  realized and publicly observed.  It changes only the institutional barrier
  in \eqref{eq:p01-institutional-wedge}; the infinite/finite wedge encodes the
  effective unavailability/availability of route \(E\).

  \item[Stage 1.] \textbf{Ex ante project advancement.}  Developer \(i\)
  observes \((a_i,k_i,M)\), anticipates the market-consistent service price
  \(p_m^*\) and future route opportunities, and chooses \(x_i\).  Equation
  \eqref{eq:p01-planned-intensity} determines the expected measure of viable
  projects reaching the route-planning stage.

  \item[Stage 2.] \textbf{Project characteristics.}  Each planning-stage
  project draws \((q,m)\) according to \eqref{eq:p01-project-draw}.  The
  optional label \(g\) records an empirical class but does not create another
  advancement choice.

  \item[Stage 3.] \textbf{Commercialization organization.}  After observing
  \((q,m)\) and the manufacturing-service price, the developer chooses a route
  from the domain in \eqref{eq:p01-route-domain}. Actual holder--producer
  separation is observed only if the realized route is \(E\), after the
  Stage~1 choice.

  \item[Stage 4.] \textbf{Downstream realization.}  A retained project
  realizes downstream commercial value with the exogenous probability in
  \eqref{eq:p01-downstream-probability}.  Planning-stage projects, route
  assignments, and realized products are distinct objects.

  \item[Stage 5.] \textbf{Manufacturing-service consistency.}  Aggregate
  entrusted-capacity demand and qualified-capacity supply must support the same
  \(p_m^*\) that developers anticipated.  This is a simultaneous equilibrium
  consistency condition, not an unanticipated event occurring after Stage~4.
\end{enumerate}

Upstream scientific research and patent generation may have occurred before
Stage~0 and may help explain predetermined \(a_i\).  They are outside the
baseline decision timing.

\subsection{Mechanism and exclusion boundaries}
\label{subsec:p01-mechanism-boundary}

The only direct reform arrow is
\begin{equation}
  M\longrightarrow \tau_E(M)
  \quad\text{whose infinite/finite value encodes route }E\text{ availability}.
  \label{eq:p01-only-direct-policy-arrow}
\end{equation}
The intended ex ante channel, conditional on the value and optimization results
derived below, is
\begin{equation}
  M
  \longrightarrow
  \text{anticipated availability/value of }E
  \longrightarrow
  \Omega_i
  \longrightarrow
  x_i^*.
  \label{eq:p01-anticipated-channel}
\end{equation}
The later realization order is
\begin{equation}
  \begin{aligned}
    x_i^*
    &\longrightarrow
    \text{planning-stage projects}
    \longrightarrow r_i^* \\
    &\longrightarrow
    \text{observed holder--producer separation}
    \longrightarrow
    \text{realized products}.
  \end{aligned}
  \label{eq:p01-realization-channel}
\end{equation}
Equations \eqref{eq:p01-anticipated-channel} and
\eqref{eq:p01-realization-channel} are timing and dependency statements, not
unconditional comparative-static propositions.

In particular, \(M\) does not directly change \(a_i\), \(k_i\), \(z_j\),
\(q\), \(m\), \(F\), \(\rho_g\), \(F_g\), \(s(q)\), or \(s_g(q)\).  It does
not change credit supply and does not directly lower \(p_m^*\). Nothing in this section implies an increase in
patent applications, upstream scientific research, or one novelty class
relative to the other.

The exact-category registry, dimensions, and forbidden interpretations are
fixed by the definitions above.  The equilibrium section below closes the
deferred fixed point without adding another policy channel.
